\section{Comparative Orderings of the Angelic Orders}\label{app:orderings}

Scripture names the orders in several places without setting all nine in
one list. The Fathers gather the names; Dionysius and Gregory the Great
rank them; the later Fathers, the medieval doctors, Aquinas, and Dante
follow one ranking or the other. The first table gives each witness's list
with its locus, highest order first where the witness ranks the orders and
in the witness's own sequence where he does not; the second and third set
the ranks side by side. Parker's English of the \emph{Celestial Hierarchy}
calls the Dominations ``Lordships,'' the Virtues ``Powers,'' and the Powers
``Authorities''; Salmond's Damascene calls the Principalities ``Rulers.''

\begingroup
\footnotesize
\setlength{\LTpre}{0.6\baselineskip}
\begin{longtable}{@{}>{\raggedright\arraybackslash}p{0.15\linewidth}>{\raggedright\arraybackslash}p{0.17\linewidth}>{\raggedright\arraybackslash}p{0.40\linewidth}>{\raggedright\arraybackslash}p{0.21\linewidth}@{}}
\toprule
\textbf{Witness} & \textbf{Locus} & \textbf{Orders} & \textbf{Arrangement}\\
\midrule
\endfirsthead
\toprule
\textbf{Witness} & \textbf{Locus} & \textbf{Orders} & \textbf{Arrangement}\\
\midrule
\endhead
\bottomrule
\endfoot
St Paul & Eph 1:21 & ``principality and power and virtue and dominion'' & Four middle orders, named from the lowest upward (\ST{I q.~108 a.~6})\\
St Paul & Col 1:16 & ``thrones, or dominations, or principalities, or powers'' & Four orders, named from the highest downward (\ST{I q.~108 a.~6})\\
Cyril of Jerusalem & \emph{Myst.\ Cat.}\ V.6 & Angels, Archangels, Virtues, Dominions, Principalities, Powers, Thrones; the Cherubim; the Seraphim & Named in the eucharistic preface; not ranked\\
Gregory Nazianzen & \emph{Or.}\ 28.31 & Angels and Archangels, Thrones, Dominions, Princedoms, Powers, Splendours, Ascents, Intelligent Powers & Not ranked; three names beyond the nine\\
Ambrose & \emph{De fide} V.13.166 & angels, archangels, cherubim and seraphim, thrones and dominions and powers & Not ranked\\
Dionysius & \emph{CH}~6.2; 7.1; 8.1; 9.1--2 & Seraphim, Cherubim, Thrones; Lordships, Powers, Authorities; Principalities, Archangels, Angels & Three triads, ranked\\
Gregory the Great & \emph{Hom.\ in Ev.} 34.7, 10 & Seraphim, Cherubim, Thrones, Dominations, Principalities, Powers, Virtues, Archangels, Angels & Ranked; written from the Angels upward\\
Gregory the Great & \emph{Moralia} XXXII.23.48 & Angels, Archangels, Thrones, Dominations, Virtues, Princedoms, Powers, Cherubim, Seraphim & The nine stones of Ezek 28:13; not ranked\\
Isidore & \emph{Etym.}\ VII.5.4; 17--24 & \latin{angeli, archangeli, throni, dominationes, virtutes, principatus, potestates, cherubim et seraphim} (5.4) & List not ranked; the exposition (5.17--24) ascends in Gregory's sequence\\
John Damascene & \emph{De fide orth.}\ II.3 & Seraphim, Cherubim, Thrones; Dominions, Powers, Authorities; Rulers, Archangels, Angels & Dionysius's triads, named from him\\
Hugh of St Victor & \emph{De sacr.}\ I.5.30 & Seraphim, Cherubim, Thrones, Dominations, Virtues, Powers, Principalities, Archangels, Angels & Written from the Angels upward; Dionysius's sequence\\
Bernard & \emph{De consid.}\ V.4.8 & Seraphim, Cherubim, Thrones, Dominions, Principalities, Powers, Virtues, Archangels, Angels & Expounded from the Angels upward; Gregory's sequence (also \emph{In Cant.}\ 19.2--5)\\
Bernard & \emph{In Cant.}\ 27.5 & Angels, Archangels, Virtues, Dominations, Principalities, Powers, Thrones; Cherubim; Seraphim & Not ranked\\
Hildegard & \emph{Scivias} I.6 & Seraphim, Cherubim; Thrones, Dominations, Principalities, Powers, Virtues; Archangels, Angels & Concentric choirs, two about five about two, from the outer choir inward; Gregory's sequence\\
Peter Lombard & \emph{Sent.}\ II d.~9 cc.~1--2 & Seraphim, Cherubim, Thrones; Dominations, Principalities, Powers; Virtues, Archangels, Angels & Triads ascribed to Dionysius in Gregory's placing; first list (c.~1) written as Hugh's\\
Bonaventure & \emph{Brev.}\ II.8.1, 3 & Thrones, Cherubim, Seraphim; Dominations, Virtues, Powers; Principalities, Archangels, Angels & Dionysius's triads; the first triad named from the Thrones upward\\
Thomas Aquinas & \ST{I q.~108 a.~6} & Seraphim, Cherubim, Thrones; Dominations, Virtues, Powers; Principalities, Archangels, Angels & Dionysius's order (\emph{sed contra}); Gregory's held reasonable (co.)\\
Dante & \emph{Par.}\ XXVIII.98--126 & Seraphim, Cherubim, Thrones; Dominions, Virtues, Powers; Principalities, Archangels, Angels & Dionysius's order (Longfellow's English)\\
Chaplet of St Michael & \emph{Raccolta} (1910) n.~291, pp.~261--263 & Seraphim, Cherubim, Thrones, Dominations, Powers, Virtues, Principalities, Archangels, Angels & Nine salutations, highest first; a third placing in the middle (Powers 5, Virtues 6, Principalities 7)\\
\end{longtable}
\endgroup

\begin{center}
\footnotesize
\textbf{The Dionysian line} (1 = highest)\\[0.4\baselineskip]
\begin{tabular}{@{}lcccc@{}}
\toprule
\textbf{Order} & \textbf{Dionysius} & \textbf{Damascene} & \textbf{Aquinas} & \textbf{Dante}\\
 & \emph{CH}~7--9 & II.3 & q.~108 a.~6 s.c. & \emph{Par.}\ XXVIII\\
\midrule
Seraphim & 1 & 1 & 1 & 1\\
Cherubim & 2 & 2 & 2 & 2\\
Thrones & 3 & 3 & 3 & 3\\
Dominations & 4 & 4 & 4 & 4\\
Virtues & 5 & 5 & 5 & 5\\
Powers & 6 & 6 & 6 & 6\\
Principalities & 7 & 7 & 7 & 7\\
Archangels & 8 & 8 & 8 & 8\\
Angels & 9 & 9 & 9 & 9\\
\bottomrule
\end{tabular}
\end{center}

\begin{center}
\footnotesize
\textbf{The Gregorian line} (1 = highest)\\[0.4\baselineskip]
\begin{tabular}{@{}lcccc@{}}
\toprule
\textbf{Order} & \textbf{Gregory} & \textbf{Bernard} & \textbf{Hildegard} & \textbf{Lombard}\\
 & \emph{Hom.\ in Ev.}\ 34.7, 10 & \emph{De consid.}\ V.4.8 & \emph{Scivias} I.6 & II d.~9 cc.~1--2\\
\midrule
Seraphim & 1 & 1 & 1 & 1\\
Cherubim & 2 & 2 & 2 & 2\\
Thrones & 3 & 3 & 3 & 3\\
Dominations & 4 & 4 & 4 & 4\\
Principalities & 5 & 5 & 5 & 5\\
Powers & 6 & 6 & 6 & 6\\
Virtues & 7 & 7 & 7 & 7\\
Archangels & 8 & 8 & 8 & 8\\
Angels & 9 & 9 & 9 & 9\\
\bottomrule
\end{tabular}
\end{center}

Hugh of Saint Victor's list and Bonaventure's triads follow the Dionysian
line (\emph{De sacr.}\ I.5.30; \emph{Brev.}\ II.8.1); Isidore's exposition
and Bernard's sermon on the love of the orders follow the Gregorian
(\emph{Etym.}\ VII.5.17--24; \emph{In Cant.}\ 19). Hildegard's column
counts the choirs from the outermost inward. The Lombard's triads, which he
ascribes to Dionysius, keep Gregory's middle (\emph{Sent.}\ II d.~9 c.~1).
The chaplet of Saint Michael, indulgenced by Pius~IX in 1851, salutes the
choirs from the Seraphim down and sets the Powers fifth, the Virtues
sixth, and the Principalities seventh, a sequence that is neither
Dionysius's nor Gregory's (\emph{Raccolta}, 1910, n.~291;
\S\ref{sec:devotion}).

Dionysius and Gregory agree on seven of the nine places. Both set the
Seraphim, Cherubim, and Thrones first, the Dominations fourth, the Powers
sixth, and the Archangels and Angels last. They differ on the Virtues and
the Principalities: Dionysius sets the Virtues fifth, next beneath the
Dominations, and the Principalities seventh, at the head of the lowest
triad (\emph{CH}~8.1, 9.1); Gregory sets the Principalities fifth and the
Virtues seventh (\emph{Hom.\ in Ev.} 34.7, 10). Aquinas receives both.
``Each of these placings may claim authority from the words of the
Apostle'': in Ephesians, which ``enumerates the middle orders, beginning
from the lowest,'' the Virtues stand between the Powers and the
Dominations, as in Dionysius; in Colossians, ``numbering the same orders
from the highest,'' the Principalities stand between the Dominations and the
Powers, as in Gregory (\ST{I q.~108 a.~6}). He gives the reason of each. Dionysius
orders the names by their relation to God, to the universal disposition of
what is to be done, and to its execution, so that the Principalities stand
first among those who execute; Gregory orders them by the things the
Dominations dispose, so that the Principalities, ``who rule even over good
spirits,'' come next beneath the Dominations, then the Powers, ``who coerce
the evil spirits,'' then the Virtues, ``which have power over corporeal
nature in the working of miracles'' (\ST{I q.~108 a.~6}). And he concludes
that ``little or no difference exists in reality between the dispositions
of the orders according to Dionysius and Gregory,'' since Gregory's
Principalities answer to Dionysius's Virtues and Gregory's Virtues to
Dionysius's Principalities (\ST{I q.~108 a.~6 ad~4}). Dante gives the same
judgment in verse. Beatrice names the nine as Dionysius named them, and
adds:

\begin{quote}
But Gregory afterwards dissented from him;\\
\hspace*{1.2em}Wherefore, as soon as he unclosed his eyes\\
\hspace*{1.2em}Within this heaven, he at himself did smile.
(\emph{Par.}\ XXVIII.133--135)
\end{quote}

\noindent See \S\ref{sec:q108} and Appendix~\ref{app:orders} on the orders and
their offices, and \S\S\ref{sec:ch7}--\ref{sec:ch9} on the Dionysian chapters.
